\documentclass[11pt]{article}
\usepackage[T1]{fontenc}
\usepackage[utf8]{inputenc}
\usepackage{lmodern}
\usepackage[margin=1in]{geometry}
\usepackage{amsmath,amssymb,amsthm,booktabs,microtype}
\usepackage[hidelinks]{hyperref}
\hypersetup{pdftitle={A finite-width enumeration of ascent sequences avoiding 0120 and 0121},
  pdfauthor={Xiaojian Zeng}}
\newtheorem{theorem}{Theorem}
\newtheorem{lemma}[theorem]{Lemma}
\newtheorem{proposition}[theorem]{Proposition}
\theoremstyle{remark}
\newtheorem{remark}[theorem]{Remark}
\DeclareMathOperator{\asc}{asc}
\newcommand{\A}{\mathcal A}
\newcommand{\D}{\mathcal D}
\newcommand{\one}{\mathbf 1}
\title{A finite-width enumeration of ascent sequences\\avoiding 0120 and 0121}
\author{Xiaojian Zeng\\
  \small PBC School of Finance, Tsinghua University, 100190\\
  \small\href{mailto:xiaojianzeng@foxmail.com}{xiaojianzeng@foxmail.com}}
\date{}
\begin{document}
\maketitle

\begin{abstract}
We give an all-length structural enumeration of ascent sequences that
simultaneously avoid the classical patterns 0120 and 0121. After the first
entry greater than one, a sufficient continuation state records the number
of eligible letters and the ranks of those already seen. Its width never
increases. At exit width $K$, the resulting finite system has
$2^K-K-1$ states and counts suffixes of arbitrary length. Decomposition at
the maximal binary prefix gives a locally finite matrix generating series
and a finite binomial--matrix formula for every coefficient. All
coefficients through degree $N$ can be computed in
$O(NL2^L+N^3)$ arithmetic operations, where
$L=\lfloor(N+1)/2\rfloor$, without generating the words. The result covers
the full avoidance class, including sequences that are not restricted
growth functions. The generating series is formal; no scalar rationality,
algebraicity, or analytic convergence assertion is made.
\end{abstract}

\section{The class and its enumeration}

For a finite word $w=w_1\cdots w_m$, let
$\asc(w)=|\{i:1\leq i<m,\ w_i<w_{i+1}\}|$.
An \emph{ascent sequence} is a word of nonnegative integers with $w_1=0$
and $w_i\leq 1+\asc(w_1\cdots w_{i-1})$ for $i\geq2$.
These sequences were introduced in connection with $(2+2)$-free posets
by Bousquet-M\'elou, Claesson, Dukes, and Kitaev~\cite{BMC2010};
pattern avoidance in them was studied by Duncan and
Steingr\'imsson~\cite{DS2011}.
Pattern occurrences here are subsequences, with order and equality
preserved and with no adjacency requirement.

Let $\A_n$ be the ascent sequences of length $n$ avoiding both 0120 and
0121. Equivalently, their entries satisfy
\begin{equation}\label{eq:avoid}
 \nexists\ i<j<k<\ell
 \quad\text{such that}\quad
 w_i<w_j<w_k\quad\text{and}\quad w_\ell\in\{w_i,w_j\}.
\end{equation}
Include $\A_0=\{\epsilon\}$, put $a_n=|\A_n|$, and write
$F(z)=\sum_{n\geq0}a_nz^n$.
Liu, Kitaev, and Zhang~\cite[Table~1 and Section~6]{LKZ2026}
list the initial values of this class and leave its enumeration
unresolved in their April 2026 preprint. The formulas below provide a
uniform structural enumeration for this exact class.

Generating trees organize recursive constructions by succession
rules~\cite{Banderier2002}; finite transition matrices then enumerate
paths by the transfer-matrix method~\cite[Section~4.7]{Stanley2012}.
The class-specific step here is a sufficient state of eligible ranks
whose width cannot increase after the first nonbinary entry. At every
fixed exit width, this gives a finite system for suffixes of all lengths.

For an integer $K\geq2$, define
\[
 \D_K=\{(k,H):2\leq k\leq K,\quad
       \varnothing\ne H\subseteq\{0,\ldots,k-2\}\}.
\]
For $s=(k,H)$, put $m=\min H$. For each $t\in\{0,\ldots,k-1\}$, set
\begin{equation}\label{eq:transition}
 f(s,t)=
 \begin{cases}
 (k,H\cup\{t\}),&t\leq m,\\[2pt]
 \bigl(k+1-r,\{h-r:h\in H,\ h\geq t\}\cup\{t-r\}\bigr),&t>m,
 \end{cases}
\end{equation}
where $r=|\{h\in H:h<t\}|$ in the second case.
The rule takes values in $\D_K$, as proved below. Let $T_K$ be the
matrix with rows and columns indexed by $\D_K$ and entries
\[
 T_K(s,u)=|\{t:0\leq t<k,\ f(s,t)=u\}| \qquad(s=(k,H)).
\]
Thus repeated successors retain their multiplicity. Let $b_K$ be the row
whose entries are one at $(K,\{h\})$, $0\leq h\leq K-2$, and zero
elsewhere, and let $\one$ be the all-ones column of the appropriate size.

\begin{theorem}\label{thm:main}
In $\mathbb Z[[z]]$ one has
\begin{equation}\label{eq:series}
 F(z)=\frac{1-z}{1-2z}
   +\sum_{K\geq2}\left(\frac{z}{1-z}\right)^{2K-1}
      b_K(I-zT_K)^{-1}\one.
\end{equation}
Equivalently, $a_0=1$ and, for every $n\geq1$,
\begin{equation}\label{eq:coefficients}
 a_n=2^{n-1}+
 \sum_{K=2}^{\lfloor(n+1)/2\rfloor}
 \sum_{d=0}^{n-2K+1}
 \binom{n-d-1}{2K-2}\,b_KT_K^d\one.
\end{equation}
The matrix $T_K$ has $2^K-K-1$ states. Its continuation functions are
rational, and only $K\leq\lfloor(N+1)/2\rfloor$ contributes through
degree $N$ in~\eqref{eq:series}.
\end{theorem}

The formula is an infinite, locally finite sum of explicitly specified
finite matrix expressions. The theorem makes no assertion that the full
series $F$ is a rational or algebraic function.

\section{Eligible letters and the quotient state}

For a nonempty avoiding prefix $w$, write $a=\asc(w)$, let $S$ be its set
of letters, and let $\lambda$ be its final letter. Define the banned set
and eligible set by
\begin{equation}\label{eq:eligible}
 B(w)=\bigcup_{\substack{i<j<k\\w_i<w_j<w_k}}\{w_i,w_j\},
 \qquad E(w)=\{0,\ldots,a+1\}\setminus B(w).
\end{equation}
Since an occurrence newly created by an append must end at that append,
\eqref{eq:avoid} gives the exact criterion
\begin{equation}\label{eq:append}
 wv\text{ is an avoiding ascent sequence}\quad\Longleftrightarrow
 \quad v\in E(w).
\end{equation}
Call a prefix \emph{postbinary} if it contains a letter at least two.

\begin{lemma}\label{lem:invariants}
Every nonempty ascent sequence satisfies $\max(w)\leq\asc(w)$. For a postbinary
ascent sequence,
\begin{equation}\label{eq:banned}
 B(w)=\{0\}\cup
 \{u>0:\text{an occurrence of $u$ has a larger letter later in $w$}\}.
\end{equation}
If the postbinary sequence avoids 0120 and 0121, then
\begin{equation}\label{eq:last}
 \lambda=\min(S\cap E(w)),\qquad a+1\in E(w)\setminus S.
\end{equation}
\end{lemma}

\begin{proof}
The maximum bound holds for $w=0$. Suppose it holds before appending
$v\leq a+1$. If $v\leq a$, the old bound persists. If $v=a+1$, then $v$
exceeds the old maximum, so the append creates an ascent and equals the
new ascent count. This proves the bound by induction.

The first nonzero letter of an ascent sequence must be one. Thus, before
the first letter at least two, both zero and one have appeared in this
order, giving a triple that bans zero. If a positive letter $u$ has a
larger letter $v$ later, the initial zero precedes its occurrence and
forms an increasing triple $(0,u,v)$. Conversely, either positive lower
value in an increasing triple has a larger letter later. This proves
\eqref{eq:banned}. Binary prefixes have $B(w)=\varnothing$.

The final letter $\lambda$ of an avoiding prefix cannot belong to $B(w)$:
any triple banning it has third value larger than $\lambda$, hence ends
before the final position, which would complete a forbidden occurrence.
In the postbinary case, $\lambda>0$, and every smaller seen value is banned
by~\eqref{eq:banned}, using the final $\lambda$ as the larger later letter.
This proves the first identity in~\eqref{eq:last}. The maximum bound
shows that $a+1$ is unseen; since banned letters have occurred, it is
also eligible.
\end{proof}

Rank the elements of $E(w)$ increasingly, starting at rank zero. Set
$k=|E(w)|$ and let $H$ be the ranks of the letters in $S\cap E(w)$.
Lemma~\ref{lem:invariants} shows that $k\geq2$, that
$\varnothing\ne H\subseteq\{0,\ldots,k-2\}$, and that the rank of
$\lambda$ is $\min H$. Denote this label by $s(w)=(k,H)$.

\begin{lemma}\label{lem:quotient}
Appending the eligible letter of rank $t$ changes $s(w)$ to
$f(s(w),t)$ as in~\eqref{eq:transition}. The width $k$ never increases.
For every $s\in\D_K$, the number of legal suffixes of length $d$ from
a prefix with label $s$ is $(T_K^d\one)_s$.
\end{lemma}

\begin{proof}
For the appended numerical value $v$, equation~\eqref{eq:banned} gives
\begin{equation}\label{eq:update}
 B(wv)=B(w)\cup\{u\in S:u<v\}.
\end{equation}
Zero is already banned, so it causes no exception in this formula.

If $t\leq\min H$, then $v\leq\lambda$. There is no ascent, and no seen
eligible value lies below $v$. Thus $E(w)$ is unchanged and the seen
ranks become $H\cup\{t\}$.

If $t>\min H$, the append is an ascent. Exactly
$r=|\{h\in H:h<t\}|$ seen eligible values are deleted by~\eqref{eq:update}.
The increase of the ascent bound adds one new unseen top value, $a+2$.
The width therefore becomes $k+1-r$. Each surviving seen rank $h\geq t$
shifts down by $r$, and the appended value has rank $t-r$. This is the
second branch of~\eqref{eq:transition}. Since $r\geq1$, the width cannot
increase.

The domain is closed even for abstract labels: in the second branch,
$r\leq t\leq k-1$, so $k+1-r\geq2$, and all seen ranks are strictly
below the new top rank $k-r$. In the first branch,
$t\leq\min H\leq k-2$, giving closure immediately.

Each rank represents exactly one eligible numerical choice. Consequently
the rule describes the successor \emph{multiset}, even when several
choices lead to the same label. Induction on $d$, starting with one
empty suffix, gives $(T_K^d\one)_s$. No information about numerical gaps
is assumed irrelevant: its effect on the ascent bound has been accounted
for explicitly in the transition calculation.
\end{proof}

Every label $(K,H)$ is realizable. Start with $(01)^{K-1}$, append
$\max(H)+2$, and then append the remaining values $h+2$ in decreasing
order. After the first of these appends the eligible set is
$\{2,\ldots,K+1\}$; the subsequent appends are nonascending and leave
it unchanged. This realizes precisely $H$. Applying the construction
with $k$ in place of $K$ realizes every label of $\D_K$.

\section{Binary initialization and assembly}

\begin{proof}[Proof of Theorem~\ref{thm:main}]
Every binary word starting in zero is an ascent sequence and avoids both
patterns. These words contribute $2^{n-1}$ at length $n\geq1$ and,
together with the empty word, give $(1-z)/(1-2z)$.

Each remaining word has a unique maximal binary prefix $p$, followed by
its first nonbinary value $v$ and then a suffix. Suppose $\asc(p)=a$.
An exit is possible exactly when $a\geq1$ and $2\leq v\leq a+1$.
It creates an ascent and bans exactly zero and one. The eligible set of
$pv$ is $\{2,\ldots,a+2\}$, with only $v$ seen. Thus
\begin{equation}\label{eq:initial}
 s(pv)=(a+1,\{v-2\}).
\end{equation}
At $K=a+1$, the exit choices give each of the $K-1$ singleton labels
in $b_K$ once.

Put $q=z/(1-z)$, the generating series for a nonempty constant run.
A binary prefix starting in zero with $a\geq1$ ascents has $2a$ runs
if it ends in one and $2a+1$ runs if it ends in zero. Its combined
weight is therefore
\[
 B_a(z)=q^{2a}+q^{2a+1}=\frac{z^{2a}}{(1-z)^{2a+1}}.
\]
The first nonbinary letter contributes a further $z$, so
$zB_{K-1}(z)=q^{2K-1}$. Lemma~\ref{lem:quotient} gives the suffix
column
\[
 \sum_{d\geq0}z^dT_K^d\one=(I-zT_K)^{-1}\one.
\]
Multiplying the prefix, exit, and suffix weights proves~\eqref{eq:series}.
The decomposition is unique, so no word is counted twice or omitted.

There are
\[
 |\D_K|=\sum_{k=2}^K(2^{k-1}-1)=2^K-K-1
\]
states. The inverse $ (I-zT_K)^{-1}$ exists formally because its matrix
has identity constant term, and its entries are rational functions by
the adjugate formula. The minimum length at exit width $K$ is $2K-1$,
attained by $(01)^{K-1}$ followed by an allowed exit. Hence each summand
has order at least $2K-1$, proving local finiteness and the cutoff.
Finally,
\[
 [z^{n-d}]\left(\frac{z}{1-z}\right)^{2K-1}
 =\binom{n-d-1}{2K-2}\qquad(n-d\geq2K-1),
\]
which yields~\eqref{eq:coefficients}. All these operations are in formal
power series; no analytic convergence is used.
\end{proof}

\begin{proposition}\label{prop:cost}
For $N\geq3$ and $L=\lfloor(N+1)/2\rfloor$, all coefficients through
degree $N$ can be computed in $O(NL2^L+N^3)$ arithmetic operations, with
$O(2^L+N)$ live counting values apart from stored transitions.
\end{proposition}

\begin{proof}
Use the single domain $\D_L$. Each state has at most $L$ outgoing
choices. Building the transitions by scanning rank sets costs
$O(L^22^L)$ elementary work, absorbed by $O(NL2^L)$.
At each length, propagate postbinary counts along these transitions.
This takes $O(NL2^L)$ additions over $N$ levels.
Keep binary counts separately by $(a,\lambda)$, where
$\lambda\in\{0,1\}$. There are $O(N)$ such labels per level, each with
two binary extensions and $O(N)$ exits specified by~\eqref{eq:initial}.
Their total cost is $O(N^3)$. Only the present and next counting arrays
are needed. This proves the claims. The bound counts arithmetic
operations, not bit operations on bounded-size integers.
\end{proof}

\section{Small systems and exact checks}

For $K=2$, the only label is $(2,\{0\})$, and both choices return to it.
Thus $b_2(I-zT_2)^{-1}\one=1/(1-2z)$.
For $K=3$, let $c_0,c_1,c_{01}$ be the continuation series of
$(3,\{0\})$, $(3,\{1\})$, and $(3,\{0,1\})$, respectively, and put
$c_2=1/(1-2z)$. The transition rule gives
\begin{align*}
 c_0&=1+z(2c_0+c_1),&
 c_1&=1+z(2c_1+c_{01}),&
 c_{01}&=1+z(c_{01}+c_0+c_2).
\end{align*}
In particular,
\[
 b_3(I-zT_3)^{-1}\one=c_0+c_1
 =\frac{2-8z+12z^2-7z^3}{(1-2z)(1-5z+8z^2-5z^3)}.
\]
This example exhibits both multiplicities and transitions to smaller width.
The initial coefficients are shown in Table~\ref{tab:counts}.

\begin{table}[ht]
\centering
\begin{tabular}{r r@{\qquad}r r}
\toprule
$n$&$a_n$&$n$&$a_n$\\
\midrule
0&1&9&1980\\
1&1&10&5945\\
2&2&11&18470\\
3&5&12&59306\\
4&13&13&196547\\
5&34&14&671306\\
6&90&15&2359338\\
7&244&16&8519890\\
8&683&&\\
\bottomrule
\end{tabular}
\caption{Coefficients of the full simultaneous-avoidance class.}
\label{tab:counts}
\end{table}

The exact checks reconstruct the objects from their definitions and use
separate implementations, a verification design also used in the finite
geometric certificates of Zeng et al.~\cite{Zeng2026}.
Here the computations test the enumeration and its implementation;
the all-length conclusion follows from the preceding induction.
The accompanying verification record has the following coverage:
\begin{itemize}
\item A literal quadruple test of every ascent sequence through length
9 includes all 31,240 sequences at length 9 and retains 1,980 of them.
A separate unmerged word census using increasing triples reaches
length 13 and agrees with~\eqref{eq:coefficients}.
\item Successor multisets and eligible-set invariants were checked on
all 8,997 avoiding prefixes of lengths 1 through 10, covering 27,466
legal edges. Closure and realizability were checked for all 247 labels
through width 8.
\item A reference recurrence retaining increasing value pairs checks
all reachable exact states through prefix length 15 and yields counts
through length 16. Its last level has 16,575 states representing
2,359,338 words, with 67,433 outgoing edges checked.
\item Separate binomial--matrix and mixed binary/rank computations
agree through degree 28, where $a_{28}=234799242479776$.
Exact rational systems were checked at widths 2, 3, and 4.
\end{itemize}
The larger exact-state computation merges equivalent reference states;
it is not a literal census of all words at those lengths. Structured
boundary checks through width 32 and 40 seeded walks of 160 steps each
provide additional finite stress tests, with no exhaustive claim for
those larger ranges.

The ancillary \texttt{coefficients.py} is a short standard-library
implementation of the mixed binary/rank recurrence. The detailed
reference programs and their coverage reports are supplied separately.
The symbolic checks used Python 3.10.12 and SymPy 1.14.0;
SymPy is described by Meurer et al.~\cite{SymPy2017}.
All reported count comparisons use integers, and the rational residual
checks are exact. There are no floating-point fits in the enumeration.

\section{Scope and comparison limits}

The eligible-rank quotient and the nonincreasing width give an
all-length structural answer to the class left unenumerated in the
checked April preprint~\cite{LKZ2026}. The finite matrix for a fixed
exit width handles arbitrarily long suffixes, while the coefficient
algorithm uses substantially fewer states than the number of words:
already the binary members alone number $2^{N-1}$ at length $N$.
No minimality assertion for the quotient is intended.

The full class cannot be replaced by restricted growth functions
(RGFs), for which each first occurrence of a positive letter has an
earlier occurrence of its predecessor. The member $01013$ is not an
RGF. Accordingly, the shorter-pattern reductions of
\cite[Lemma~2.2]{LKZ2026}, which require the RGF hypothesis, do not
apply to the entire present class. Mansour and Shattuck's inspected
preprint~\cite{MS2025} studies simultaneous avoidance of 021 and a
length-four pattern. The word $010132$ belongs to our class but
contains 021 at positions $(1,5,6)$ and 0132 at $(1,2,5,6)$.
Thus either additional avoidance condition would discard valid inputs.

The literature comparison is bounded by the primary material inspected
on 13 September 2026. No equivalent enumeration was found in that
material. The full text of Mansour's later study of 0132 avoidance
\cite{Mansour2026} was unavailable for comparison, and changes between
the inspected 021 preprint and its journal version were not checked.
These limits prevent a claim of established worldwide priority or a
complete exclusion of overlap with those works. The self-contained
enumeration does not depend on a theorem from either work. A simpler
scalar expression for $F$, its algebraic classification, and its
analytic growth are not determined here.

\paragraph{Computational and AI assistance.}
OpenAI Codex assisted with the literature search, proof development,
exact-check programs, separate model-context mathematical checking,
and manuscript preparation. A separate model context is not external
peer review or statistically independent evidence. The theorem is
supported by the written proof; the explicitly delimited computations
check finite instances and implementations. No proof assistant was used.

\bibliographystyle{plainurl}
\bibliography{references}
\end{document}
